% Generated draft: requires mathematical and source review.
\documentclass[11pt]{article}
\usepackage[margin=25mm]{geometry}
\usepackage{fontspec}
\setmainfont{Latin Modern Roman}
\usepackage{amsmath,amssymb,mathtools,unicode-math}
\setmathfont{Latin Modern Math}
\usepackage{enumitem,needspace,xurl,hyperref}
\setlength{\parindent}{0pt}
\setlength{\parskip}{5pt}
\setlength{\emergencystretch}{3em}
\newcommand{\N}{\mathbb{N}}
\newcommand{\Z}{\mathbb{Z}}
\newcommand{\Q}{\mathbb{Q}}
\newcommand{\R}{\mathbb{R}}
\newcommand{\C}{\mathbb{C}}
\newcommand{\F}{\mathbb{F}}
\newcommand{\A}{\mathbb{A}}
\newcommand{\PP}{\mathbb{P}}
\newcommand{\E}{\mathbb{E}}
\newcommand{\eps}{\varepsilon}
\newcommand{\dd}{\,\mathrm{d}}
\begin{document}
\section*{UnsolvedMath ID: 30004206}
\textbf{Source problem code:} catalogue entry
\subsection{Definitions and mathematical statement}
Let \(k\) be a field of characteristic zero. Write \(\rho\) for the class denoted by \(\rho\) in the supplied source. The hypothesis that \(\rho\) is nonnilpotent means
\[
\rho^n\ne 0\qquad\text{for every integer }n\ge 1
\]
in its ambient graded ring. The supplied record does not specify that ring. In the customary mod-\(2\) motivic-cohomology interpretation, \(\rho=[-1]\) belongs to \(K_1^M(k)/2\cong H^{1,1}(k;\mathbb Z/2)\), where \(K_*^M(k)\) is Milnor \(K\)-theory and \([-1]\) is the symbol of the unit \(-1\). Under this interpretation, nonnilpotence means \([-1]^n\ne0\) in \(K_n^M(k)/2\) for every \(n\ge1\). This interpretation requires confirmation from the original report.

Let \(\eta\) denote the motivic Hopf class. Inverting \(\eta\) means making multiplication by this class invertible; for a motivic spectrum \(E\), the resulting localization is denoted \(E[\eta^{-1}]\).

\textbf{Problem.} For every field \(k\) of characteristic zero satisfying the stated nonnilpotence hypothesis, determine the profile of the \(\eta\)-periodic slice spectral sequence over \(k\): determine its differentials, the extension data needed to reconstruct its abutment, and the abutment itself.

\textbf{Unresolved specifications.} The record does not identify the underlying motivic spectrum, its coefficient or completion conventions, the indexing of the spectral sequence, or its convergence conditions. In particular, it does not explicitly say whether the intended spectrum is the motivic sphere, nor whether periodicity is imposed on a slice spectral sequence or by first localizing a spectrum. These constructions must not be identified without justification. Consequently, the problem above preserves the supplied target but does not specify a unique spectral sequence. The supplied background also mentions a conjectured real-field profile without giving that conjecture.
\subsection{Short English statement}
For each characteristic-zero field where the source's class \(\rho\) has no vanishing positive power, determine all differentials, extensions, and the abutment of the intended \(\eta\)-periodic slice spectral sequence. The underlying spectrum and several conventions remain unspecified. The supplied record labels the problem open; this draft does not independently verify that status.
\subsection{Sources}
\begin{itemize}
\item Supplied local UnsolvedMath registry record, permanent ID \(30004206\), title ``Eta-Periodic Slice Spectral Sequences over General Fields,'' in the Oberwolfach Reports Open Problems collection. Exact source problem code: \texttt{catalogue entry}. The record gives proposed year \(2019\) and names no individual proposer.
\item Supplied source citation: \emph{Homotopy Theory (2020)}, Oberwolfach Reports, pp.~2183--2256, DOI \texttt{10.4171/owr/2019/36}. Supplied source URL: \url{https://doi.org/10.4171/owr/2019/36}. The record describes these reports as documenting mathematical workshops and their open-problem discussions.
\item Supplied literature source, ``original Oberwolfach report'': \url{https://doi.org/10.4171/OWR/2019/36}.
\item Supplied literature source, ``R-motivic v-one-periodic homotopy'': \url{https://arxiv.org/abs/2204.05937}.
\item Supplied literature assessment dated exactly \texttt{2026-08-22}: the profile over characteristic-zero fields with nonnilpotent \(\rho\), including the conjectured real-field profile, is described as unresolved in the cited and later literature. This is supplied dated curation derived from the OWR record, not an independently verified literature assessment. The supplied background expressly requests independent verification of time-sensitive status and cited claims.
\end{itemize}
\end{document}
